% Transcription — fonds Alexandre Grothendieck, shelfmark 156-4, batch 5 (pages 81-88).
% Established from the facsimile archives/batches/156-4/batch-05.pdf (a copy of
% 156-4.pdf fetched through the project's relay; no cover sheet in the batch, so
% PDF page k = folder page 80+k), per the skill transcribe-grothendieck. The
% folder is eighty-eight pages in five batches; this is the fifth and last,
% eight pages. Batches 1-2 were read for notation; batches 3-4 (pages 41-80)
% were being transcribed in parallel and did not exist when this batch was
% written, so page 81 is read against the facsimile of page 80 only: page 80
% ends on a corollary (G' ≼ G ≼ F ⟹ ...) and a « Démonstration » beginning
% « Supposons X ≤ F, », which page 81 continues. The axioms C5, C6, C7 are
% those of batch 4 (named on page 80); they are cited here, not restated.
%
% Pass: Opus 5.5 (claude-opus-5-5), 2026-09-26 — first pass, unchecked against the pages by a human.
%
% Pages skipped: none. Pages summarised: none. All eight leaves are his
% mathematics, fast drafting in ink, and all are transcribed. Page 81 is
% cancelled in two blocks by long oblique strokes (the top third and the
% whole « C 8 » statement at the foot), the lower half of page 86 and the
% top half of page 87 likewise; they are transcribed and the cancellation
% noted once per block.
%
% His pagination 81-88 runs at the top of the leaves and coincides with the
% archivists'. A struck number stands to the left of « 87 » on page 87 (read
% as a first « 88 », not certain). No date in his hand on these pages. His
% numbered runs: the axiom « C 8 », first stated and cancelled (p. 81,
% a)-b)), then restated as « C 8 (0) » with situations (a), (b) (p. 82),
% cited as C8(a), C8(b) (p. 83); on page 86 « C8 » is reused for a new axiom,
% « Axiome de recollement des subdivisions » (a first « C8 » struck above
% it). He cites C7 (pp. 83, 87), C5 (p. 83), C1 à C4 (p. 85), C6 and C6 (c)
% (p. 86). Page 83 marks two statements (*) and (**). Page 84's « Remarque »
% has conditions a), b), c) and conclusions (i); page 85 continues (ii),
% (iii). Pages 85-86 list four packets of axioms under circled signs, the
% inner sign not read (set $\circledast$ with a note). Cross-references in his
% hand: « Cor. prop. (p. 78) » (p. 81, the digit after 7 doubtful), « via
% prop. p. 78 » (p. 82), « corollaire p. 80 » (p. 83), « voir prop. p. 78 »
% (p. 88); pages 76-80 are in batch 4.
%
% What the batch is about. Figures F, G of a contrée carry two orders, ≤
% (sub-figure) and ≼ (a second order introduced in batch 4; ≪, « raffine »,
% is defined from both: G ≪ F iff G ≤ F' ≼ F for some F', p. 84). Pages
% 81-83 prove that ≪ is antisymmetric: the axiom C8 (p. 82), for a
% multistrate X, X ≤ Φ ≼ X ⟹ X = Φ and X ≼ Ψ ≤ X ⟹ X = Ψ, finishes the
% proof begun on page 80, and gives the corollary F ≪ G and G ≪ F ⟹ F = G,
% via C7 (completion of squares) and the corollary of page 80. He remarks
% that C8 is necessary if ≪ is to be an order, and that he has tried to do
% without disjunction axioms of the type C5. Pages 84-85: a « Remarque »
% abstracting this — a set with two orders ≤, ≼ satisfying a) unique
% completion of a square F ≥ G, F ≽ F', b) completion of G ≤ F, G ≽ G',
% c) F ≤ G ≼ F ⟹ F = G; then (i) ≪ is an order, (ii) the square of a) is
% cartesian for ≪, G' = Inf_≪(F', G), (iii) cartesian for ≥ implies
% cartesian for ≪. He then plans the axioms of contrées in packets: ≤ alone
% (C1-C4); the « formal » properties of ≤, ≼, without support or
% disjunction; gluing of subdivisions; disjunction. Page 86: the axiom C8 of
% gluing of subdivisions (F = Sup F_i, compatible subdivisions F'_i glue to a
% subdivision of F) and a corollary Subdiv(∐ F_i) ≃ ∏ Subdiv(F_i); heading
% « Axiomes de disjonction et de support », and « c'est le moment
% d'introduire les lieux »; a cancelled attempt (pp. 86-87) at minimal
% elements for ≼ and ≪. Page 87: definition — a lieu of C is an element of
% 𝔉 ∖ {∅_𝔉} minimal for ≪; a lieu is in ℳ and minimal for ≤, the converse
% not holding with the axioms so far (it will for « contrées sans bords »);
% ℒ the set of lieux; the ombre |F| = {x ∈ ℒ, x ≪ F} and the multiombre
% F̂ = {|X| : X ∈ F̄} ⊂ 𝔓(ℒ) (p. 88), |F| = ∪ |X|; the folder ends on a
% Proposition that F̂ is an admissible family of parts of ℒ (a) A° = A ∖ ∂A
% ≠ ∅, b) A ∩ B a union of members), stated without proof.
%
% Notation fixed once. As batch 2: \mathcal{M} (multistrates), \mathcal{L}
% (lieux, set of lieux, ℒ in his round script), \mathfrak{F} (the set of
% figures), \mathfrak{P} (power set). Refinement ≪ is \ll (double chevron
% throughout these pages; the triple \lll of batch 2 p. 38 does not occur).
% His second order, a < over a doubled curved bar, is set \preccurlyeq
% (reversed \succcurlyeq); his ≤, often with a doubled bar, is \leq. His
% bar over a figure (F̄, Ḡ_X: the set of multistrates X ≤ F, as p. 88 says)
% is \overline. The empty figure, written ∅ with an index 𝔉, is
% \emptyset_{\mathfrak{F}}. Vertical relations between the rows of his
% stacked inequalities are set in an array as the symbol placed between two
% rows, to be read from the upper term to the lower one (noted at the first
% occurrence, p. 81). His underlining is \emph.
%
% Hard pages: 81 (both cancelled blocks and the connective prose between
% them, the stacked diagram in the middle), 83 (a long oblique note in the
% left margin, read in part), 85 (a long oblique margin note and an
% interlinear insertion in (ii), read only in fragments). Readings flagged
% and not settled: the page number cited on pages 81 and 82 (76 or 78); the
% direction of several vertical relations on page 81; « suffisant » (p. 82);
% the phrase closing page 84; several words of pages 85 and 87.
%
\documentclass[11pt,a4paper]{article}
\input{../preamble/grothendieck}

\folder{156-4}
\batch{5}
\pages{81}{88}
\dating{1986}
\watermark{Édition de démonstration}
\foldertitle{[Chapitre] IV. Analysis situs (première mouture) : notes manuscrites (10/06/1986)}

\begin{document}

\page{81}
\note{la page continue la démonstration commencée au bas de la page 80 (« Supposons $X \leq F$, ») ; le tiers supérieur, jusqu'à « $X \ll G_X <$ », est barré de trois longs traits obliques}
\[
  \begin{array}{ccccc}
    X \ll & G & \leq & F \\
          & \geq & & \geq \\
          & G_X & \leq & X
  \end{array}
\]
\note{dans ce tableau et les suivants, le signe placé entre deux lignes est sa relation écrite verticalement, à lire du terme du haut vers le terme du bas ; ici $G \geq G_X$ et $F \geq X$}

\struck{Soit $G_X$ la sous-figure de $G$ image inverse de} \struck{\ill{}} \struck{la sous-figure $X$ de $F$. Par le Cor. prop. (p.~\uncertain{78}) $X$ raffine \ill{} : à la fois la sous-figure $X$ de $F$, et $G$ raffine $G_X$, donc $X \ll G_X <$}
\note{le chiffre après le 7 du renvoi est douteux (76 ou 78)}

\struck{Soit} Supposons $X \leq F$, \ill{} $X$ \uncertain{est} \ill{} de $F$, \ill{} $\overline{F}_X$, \struck{$G_X$}, $\overline{G}_X$ \ill{}, considérons les images inverses \add{(dans $F$, $G$, \struck{\ill{}})} \struck{\ill{} deux sous-figures de $F$, $G$, $\overline{G}$} des $\overline{F}$, $G$, $\overline{G}$. (Donc on trouve \struck{(par la prop.}
\note{deux lignes surchargées et barrées ; l'ordre des insertions n'est pas sûr}

\struck{$X \ll \overline{G}_X$}

\struck{\ill{}} \ill{} en particulier : $H$ ($= \overline{G}$) \ill{} figure, $H'$ ($= \overline{G}_X$) sous-figure, $X$ \add{multistrate} \struck{figure, $H$ ($X \leq H'$) \ill{}} \ill{} de la propriété universelle de la prop.
\note{passage encadré en partie et raturé ; lu par fragments}

\[
  \begin{array}{ccc}
    G_X & \leq & X \\
        &      & \succcurlyeq \\
    \overline{G}_X & \leq & F_X \\
    \geq & & \\
    X & &
  \end{array}
\]
\note{sous $G_X$ un signe surchargé, non lu ; au pied, « $X$ est une multistrate »}

On est ramené à prouver $X \leq G_X$.

\note{en haut à droite du bloc suivant, un « ($\leq F_X \leq$) » barré}
\struck{\emph{C 8} a) Soit $X \preccurlyeq G$ \ill{} multistrate, $\Phi$ \ill{}} \struck{$X$ une multistrate, $\Phi$ une subdivision de $X$, si $X \leq \Phi$ (i.e. $X$ \ill{}) on a $X = \Phi$. b) Si $X \preccurlyeq \Psi$, $X$ multistrate, $\Psi$ figure, alors $X = \Psi$. Dans les deux cas, d'après \ill{} on trouve \ill{}}
\struck{$X \leq \overline{G}_X \leq F_X = X$, donc $X = \overline{G}_X$, d'où}
\struck{$X \preccurlyeq G_X \leq X$}
\note{tout ce bloc, de « C 8 » au bas de la page, est barré de longs traits obliques ; « C 8 » souligné, un trait vertical dans la marge gauche ; il est repris à la page 82}

\page{82}
\marginal{provisoire !}
\note{la marge porte « provisoire ! » en oblique et un double trait vertical le long de l'énoncé}
\emph{C 8} (0) \quad Considérons les situations
\[
  \begin{array}{lll}
    (a) & X \leq \Phi \preccurlyeq X & \Phi \in \mathfrak{F},\ X \in \mathcal{M} \\
    (b) & X \preccurlyeq \Psi \leq X & \Psi \in \mathfrak{F},\ X \in \mathcal{M}
  \end{array}
\]
\struck{Alors} Alors on a $X = \Phi$ et $X = \Psi$.

Achevons la dém. Appliquant (a) à
\[
  X \leq F_X \preccurlyeq X
\]
on trouve $X = F_X$, d'où, comme
\[
  X \leq \overline{G}_X \leq F_X = X ,
\]
$X = \overline{G}_X$, d'où
\[
  \overline{G}_X = X \preccurlyeq G_X \leq X
\]
\note{un « $X =$ » barré devant $\overline{G}_X$}
et appliquant (b), on trouve $X = G_X$ et par suite $X \leq G_X$, cqfd.

\emph{Remarques} \quad Bien sûr, C8 est nécessaire si on veut que $\ll$ soit une relation d'ordre. Et c'est aussi \uncertain{suffisant} ;

\emph{Cor.} de (C8) \add{via prop. p.~\uncertain{78})} \quad Si $F \ll G$ et $G \ll F$, on a $F = G$.

\emph{Dém.} \quad \struck{Soit $X \in \overline{F}$ i.e. $X \in \mathcal{M}$, $X \leq F$} On va prouver $F \leq G$ ; \struck{\ill{} on aura} \struck{Il faut prouver} par symétrie $G \leq F$, donc $F = G$. \struck{pour pour $X \in \mathcal{M}$, ($X \leq F$) $\Rightarrow$ $X \leq G$.}
\note{la barre sur le $F$ de la ligne barrée est ondulée (lue comme $\overline{F}$, douteux)}
\[
  \begin{array}{ccccc}
    F & \ll & G & \ll & F \\
    \geq & & \geq & & \geq \\
    F_X & \ll & G_X & \ll & X \\
    \geq & & & & \\
    X & & & &
  \end{array}
\]
\note{ce tableau, au pied de la page, est traversé de traits verticaux et d'un trait oblique ; le premier $\ll$ est surchargé}

\page{83}
On a par définition de $\ll$
\[
  \begin{array}{ccccc}
    F & \leq & G' & \preccurlyeq & G \\
      &      &    &              & \leq \\
      &      &    &              & F' \\
      &      &    &              & \preccurlyeq \\
      &      &    &              & F
  \end{array}
\]
on complète par C7
\[
  \begin{array}{ccccc}
    F & \leq & G' & \preccurlyeq & G \\
      &      & \geq &            & \leq \\
      &      & H & \preccurlyeq & F' \\
      &      &   &              & \preccurlyeq \\
      &      &   &              & F
  \end{array}
\]
\note{le signe vertical sous $G'$ est surchargé ; lu $\geq$, douteux}

\struck{\ill{}} donc
\[
  (*) \qquad F \leq H \preccurlyeq F
\]
On va en conclure $F = H$ (grâce à C8 (a)), donc
\[
  G' = F , \quad F' = F , \quad \text{d'où}
\]
\[
  (**) \qquad F \preccurlyeq G \leq F
\]
on va en conclure $F = G$ (grâce à C8 (b)).

\marginal{Mais ça \ill{} \ill{} \ill{}, on \ill{} \ill{} \ill{} de $F$ et $H$ \ill{} $H \preccurlyeq F \leq H$ !}
\note{longue note oblique de la marge gauche, en regard de (*) et (**), précédée d'un double trait ; lue par fragments}

Prouvons d'abord ce dernier pt : il \struck{faut} suffit de prouver $F \leq G$, i.e. \struck{\ill{}} $X \in \mathcal{M}$, $X \leq F$ $\Longrightarrow$ $X \leq G$, c'est \ill{} \ill{} le corollaire p.~80.

\struck{Pour \ill{} (**)}

On a \uncertain{essayé} ici de se débrouiller aussi bien que possible, sans axiomes et raisonnements de disjonction, genre C5.

\page{84}
\emph{Remarque} \quad Soient $\mathfrak{F}$ \struck{un ens} un ensemble muni de deux relations d'ordre $\leq$, $\preccurlyeq$, satisfaisant les conditions suivantes
\begin{enumerate}
  \item[a)] Pour
  \[
    \begin{array}{ccc}
      F & \geq & G \\
      \succcurlyeq & & \\
      F' & &
    \end{array}
  \]
  on peut compléter de façon unique en
  \[
    (*) \qquad
    \begin{array}{ccc}
      F & \geq & G \\
      \succcurlyeq & & \succcurlyeq \\
      F' & \geq & G'
    \end{array}
  \]
  \note{le premier $\geq$ de a) est écrit sur un autre signe}
  \item[b)] Pour
  \[
    \begin{array}{ccc}
      G & \leq & F \\
      \succcurlyeq & & \\
      G' & &
    \end{array}
  \]
  on peut compléter (de façon pas nécessairement unique) en
  \[
    \begin{array}{ccc}
      G & \leq & F \\
      \succcurlyeq & & \succcurlyeq \\
      G' & \leq & F'
    \end{array}
  \]
  \item[c)] Si $F \leq G \preccurlyeq F$ (ou ce qui revient au même $G \preccurlyeq F \leq G$), alors $F = G$.
\end{enumerate}

\emph{Alors} : \struck{(i) Posant $F \ll \overset{\mathrm{def}}{\Longleftrightarrow} \exists$}
\begin{enumerate}
  \item[(i)] Posant $G \ll F \overset{\mathrm{def}}{\Longleftrightarrow} \exists\, F'$ avec $G \leq F' \preccurlyeq F$, la relation $\ll$ est une relation d'ordre. (\ill{} \ill{}, \ill{} \ill{} suppose que a) b))
\end{enumerate}
\note{le $G$ de « $G \ll F$ » et le $F'$ qui suit sont écrits sur d'autres lettres}

\page{85}
\begin{enumerate}
  \item[(ii)] Dans le cas a), le carré (*) \add{\ill{} $\leq$, i.e. \ldots} est cartésien pour $\ll$, i.e.
  \[
    G' = \operatorname{Inf}_{\ll}(F', G)
  \]
  \note{l'insertion interlinéaire, reliée à (*) par un trait, n'est lue qu'en partie}
  \item[(iii)] \struck{Si} Si \struck{$G$}
  \[
    \begin{array}{ccc}
      F & \geq & G \\
      \geq & & \geq \\
      F' & \geq & G'
    \end{array}
  \]
  est cartésien pour $\geq$, il est cartésien pour \struck{\ill{}} $\ll$.
  \note{le $F$ et les deux signes de la première colonne et de la première ligne sont surchargés ; les relations lues ici sont douteuses}
\end{enumerate}

\marginal{Il \ill{} que \ill{} \ill{} \struck{\ill{}} des \ill{} ?? P. ex. \ill{} $\operatorname{Inf}$ \ill{} pris dans \ill{}. Il faudrait \ill{} au clair la \uncertain{définition}.}
\note{longue note oblique de la marge gauche, en regard de (ii) et (iii), soulignée de deux traits ; lue par fragments}

Dans les axiomes des contrées, il faudrait séparer
\begin{itemize}
  \item[$\circledast$] un paquet concernant $\leq$ tout seul (C1 à C4)
  \item[$\circledast$] un paquet autour des propriétés « formelles » des relations $\leq$, $\preccurlyeq$ (\uncertain{celles} qui viennent d'apparaître) sans questions de support ni de disjonction, et en utilisant le minimum sur $\leq$ tout seul. Conditions qui assurent, quand les \uncertain{sup} \uncertain{majorés} existent pour $\leq$, que ceux-ci \uncertain{commutent} à l'induction sur une subdivision.
\end{itemize}
\note{les signes de cette liste et de celle de la page 86 sont des signes entourés dont le signe intérieur n'est pas lu}

Au besoin, si on en \ill{} \ill{}. \uncertain{Je} \ill{} \ill{} \ill{} \ill{} les axiomes minimaux, donnant le paquet des propriétés formelles « évidentes » des \ill{} de vraie signification géom.

\page{86}
\begin{itemize}
  \item[$\circledast$] Axiomes de \emph{recollement} des \struck{\ill{}} subdivisions
  \item[$\circledast$] \struck{Un} Un paquet d'axiomes de disjonction
\end{itemize}

\struck{C8} \quad ---

\emph{C8 \quad Axiome de recollement des subdivisions}

Soit $F = \operatorname{Sup} F_i$, et pour tout $i$, soit $F'_i$ une subdivision de $F_i$, telle que pour $i, j$, on ait $F'_i \mid F_i \cap F_j = F'_j \mid F_i \cap F_j$, alors $F' = \operatorname{Sup} F'_i$ existe et est une subdivision de $F$.
\note{un trait vertical dans la marge le long de l'énoncé}

\struck{Soit} (Devrait suivre C6), cela complète C6 (c)

\emph{Cor.} \quad Soit \struck{$F \amalg F_i$, alors} $F = \coprod F_i$, alors
\[
  \operatorname{Subdiv}(F) \simeq \prod_i \operatorname{Subdiv}(F_i)
\]

\emph{Axiomes de disjonction et de support.}

Pour formuler les axiomes sous forme aussi forte et \uncertain{sensible} que possible, c'est le moment d'introduire les lieux.

\struck{Définition} \struck{Rappel} \struck{On appelle} Soit $F \in \mathfrak{F}$, \add{$F \neq \emptyset_{\mathfrak{F}}$} \struck{conditions} conditions équivalentes
\begin{enumerate}
  \item[a)] $F' \preccurlyeq F \Longrightarrow F' = F$ \quad i.e. $F$ est minimal dans $\mathfrak{F}$ pour $\preccurlyeq$
  \item[b)] $F' \ll F$, \add{$F' \neq \emptyset$} $\Longrightarrow F' = F$ \quad i.e. $F$ est minimal dans $\mathfrak{F}$ \add{$\smallsetminus \lbrace\emptyset\rbrace$} pour $\ll$
\end{enumerate}
et elles impliquent que $F \in \mathcal{M}$, \struck{\ill{}} et $F$ est minimal dans $\mathcal{M}$ (\ill{} \ill{} \ill{}) pour $\leq$
\note{à partir de « Soit $F \in \mathfrak{F}$ » le bas de la page est barré de trois traits obliques ; la barre continue en haut de la page 87}

\page{87}
\note{à gauche du numéro 87, un numéro biffé, lu 88 sous réserve}
\struck{\emph{Dém.} b) $\Rightarrow$ a) trivial, prouvons a) $\Rightarrow$ b). Soit $F' \ll F$ i.e. $F' \leq \overline{F} \preccurlyeq F$. Alors (a), $\overline{F} = F$, donc $F' \leq F$. Par C7, $\exists$} \struck{On a donc $F' \leq F$ \ill{} ce qui implique $F' = F$.}
\struck{Mais si} \struck{On va le prouver tantôt.}
\struck{Ceci dit, comme $\ll$ implique $\leq$, il en résulte que (si $F \neq \emptyset$) $F$ est minimal dans $\mathfrak{F} \smallsetminus \lbrace\emptyset\rbrace$.} \struck{Pour $\leq$ : si $X' \leq$}
\note{ce haut de page est encadré et barré de longs traits obliques, dans le prolongement de la page 86}

\emph{Définition} \quad On appelle \emph{lieu} \add{de $C$} un élément de $\mathfrak{F} \smallsetminus \lbrace\emptyset_{\mathfrak{F}}\rbrace$ minimal pour $\ll$, i.e. une figure \struck{$F$} telle que \struck{$F$} $F \neq \emptyset_{\mathfrak{F}}$, et $F' \ll F$, $F' \neq \emptyset_{\mathfrak{F}}$ implique $F' = F$.
\note{un trait vertical dans la marge le long de la définition}

Ainsi, on voit qu'un lieu est \ill{} $\in \mathcal{M}$, et c'est aussi un élément de $\mathcal{M}$ minimal pour $\leq$ (mais il n'est pas vrai, dans l'état actuel des choses, i.e. avec les axiomes qu'on a introduits jusqu'à présent) que l'inverse soit vrai. (Ce sera vrai cependant pour les « contrées sans bords »\ldots)

On désigne par $\mathcal{L}$ l'ensemble des lieux.

\struck{On} Pour toute $F \in \mathfrak{F}$, on appelle \emph{ombre de $F$} l'ens. des $x \in \mathcal{L}$ tels que $x \ll F$, notation $|F|$, et \emph{multiombre} l'ens

\page{88}
des parties \struck{\ill{}} de $\mathcal{L}$
\[
  \widehat{F} = \lbrace\, |X| \mid X \in \overline{F} \ \text{i.e.}\ X \in \mathcal{M},\ X \leq F \,\rbrace
\]
Ainsi on a
\[
  |F| \subset \mathcal{L} , \qquad \widehat{F} \subset \mathfrak{P}(\mathcal{L})
\]
et on a
\[
  |F| = \bigcup_{X \in \overline{F}} |X| = \bigcup_{A \in \widehat{F}} A .
\]
\note{un chapeau barré sur le $|F|$ de la deuxième formule}

Cette relation signifie que si $x \in \mathcal{L}$ est tel que $x \ll F$, $\exists\, X \in \overline{F}$ tel que $x \ll X$, \struck{\ill{}} et pour ceci, voir prop. p.~\uncertain{78}. Il y a même un plus petit tel $X$ (cor. à la prop.).

\marginal{Mais pour la prouver, on \uncertain{utilise} \ill{} \ill{} \ill{}}
\note{note oblique de la marge gauche, en regard de ce paragraphe}

\emph{Proposition} \quad La multiombre \add{$\Phi =$} $\widehat{F}$ d'une figure est une famille \emph{admissible} de parties de $\mathcal{L}$, i.e.
\begin{enumerate}
  \item[a)] $\forall$ \struck{\ill{}} $A \in \Phi$, posant $\partial A = \bigcup_{B \subset A,\ B \neq A} B$, on a $\partial A \neq A$ i.e. $A^{\circ} \overset{\mathrm{def}}{=} A \smallsetminus \partial A \neq \emptyset$
  \item[b)] $\forall\, A, B \in \Phi$, on a $A \cap B = \bigcup_{C \in \Phi,\ C \subset A \cap B} C$
\end{enumerate}
\note{la page, dernière du dossier, s'arrête sur cet énoncé ; la démonstration ne suit pas}

\end{document}
